\section{Implementation Boundaries}
\label{sec:implementation}

The implementation connects the mechanisms through several access paths.
Their differences affect both guarantees and performance attribution.
The following boundaries come from source inspection; the companion
record identifies the inspected revisions, while runtime and failure
validation remain part of the evaluation plan.

\subsection{Buffered and direct access}

An ordinary volume resolves mutable metadata and sends file data through
buffered caching and writeback. A workspace can reuse this path while
replacing metadata lookup with private-state resolution over a sealed
catalog or packed lower. A standalone snapshot reader instead provides
read-only access without constructing a private upper. Its metadata cache
is separate from payload caching; the current read-only path does not
admit configurable payload-cache or prefetch budgets.

A second workspace implementation stores current private namespace and
mapping records transactionally and uses a dedicated direct-I/O adapter.
This is the concrete namespace and truncation path described in
Section~\ref{sec:delta}. It bypasses ordinary buffered writeback and
metadata TTL caching. Its operation sizes, complete-directory enumeration,
and live ownership set have admission bounds. It currently supports
newly created workspaces, without process-crash takeover. These limits
belong to that path, rather than the logical workspace model or every mount.

The direct path authenticates an entire upper object before returning a
subrange. Thus a smaller backend verification unit need not reduce work
for a short application read. Comparing this path with a buffered volume
changes cache and fetch behavior as well as workspace representation;
an ablation must retain one access path or control those changes explicitly.

\subsection{Publication and recovery}

Snapshot publication validates uploaded objects before installing the
lower binding and acquiring reader ownership. Recovering a binding also
requires generation and lease checks: knowing a manifest's identity does
not grant authority to resume writes. One connected recovery path reuses
the same persistent storage allocation and checks its worker identity and
scoped metadata authority. It does not provide arbitrary-worker migration.

Catalog lifecycle operations encode quiescence, object verification,
retention, and atomic version transitions. Connecting these operations
to worker handoff still requires runtime orchestration. The recovery
target is a published file state; applications must checkpoint execution
state and reopen tools separately. Source-level lifecycle operations
therefore do not by themselves demonstrate end-to-end recovery.

\subsection{Object formation and layout admission}

The ordinary uploader partitions slices using fixed volume geometry,
defaulting to $g=4\MiB$ within 64 MiB chunks. Readers use the same geometry;
an incompatible configuration is rejected. The proposed 1/4/8/16 MiB
settings expose a design space, not an adaptive sizing policy.

Backend hints are optional and require both client support and compatible
storage servers. Hint admission and new verification framing have separate
enable and fleet-compatibility gates, disabled by default. The backend
checks eligible writes and persists their effective layouts, which govern
later reads and repair. Changing admission settings does not re-encode
existing objects.

Uploader coverage matters. The ordinary uploader propagates hints; direct
workspace writes and the current snapshot producer use default settings.
Extending them must preserve conditional object creation and rejection of
conflicting writes. Direct workspace uploads also form objects from write
fragments, rather than ordinary buffered blocks, and cap those fragments
at the automatic policy's no-hint threshold.

That policy omits hints for $L_o\leq1\MiB$ and requests $B=1\MiB$
otherwise; the inspected backend already defaults to $B=1\MiB$.
Enabling hints can therefore leave the stored layout unchanged while
changing client fetch and cache policy. Experiments must record actual
$(g,L_o,B,q)$ and physical work, and separate those client effects from
backend geometry. These integration restrictions constrain current
comparisons without making any of the three mechanisms mandatory.
